\section{Diversity and Comprehensiveness Evaluation on DBLP}
\label{sec:dblp-diversity}

A benchmark is only as informative as the space of behaviours it exercises. We therefore
evaluate \textsc{QueryBloom} along \emph{three orthogonal axes of diversity} --- the
diversity of the natural-language \emph{questions} (input space), the diversity of the
\emph{answers} returned by the queries (output space), and the diversity of the
\emph{SPARQL programs} that connect them (program space) --- and, separately, along a
\emph{comprehensiveness} axis that measures how much of the knowledge graph and of the
SPARQL~1.1 language the benchmark actually touches. Diversity without coverage is
redundancy in disguise; coverage without diversity is a narrow probe repeated many times.
Both are needed, and we report both.

All results below compare \textsc{QueryBloom} against the human-curated DBLP baseline. To
remove any size confound, the larger set is uniformly downsampled so that both benchmarks
contain exactly $N{=}83$ question--query pairs. Every SPARQL statistic is computed from the
\texttt{rdflib} abstract syntax tree of the query rather than from surface text, and every
answer statistic is obtained by executing the query against the live DBLP endpoint.

\subsection{Axis 1: Semantic diversity of the questions}

We embed each NL question with \textsc{bge-large-en-v1.5} and $\ell_2$-normalise the
resulting vectors, so that inner products are cosine similarities. Over the set of
embeddings $E=\{e_1,\dots,e_N\}$ we report two complementary statistics:

\paragraph{Average pairwise distance (APD).} The mean cosine distance between all distinct
pairs,
\[
\mathrm{APD}(E)=\frac{2}{N(N-1)}\sum_{i<j}\bigl(1-e_i^{\top}e_j\bigr),
\]
a measure of the \emph{global spread} of the question set.

\paragraph{Vendi Score.} Let $K$ be the $N\times N$ cosine-similarity matrix and
$\lambda_1,\dots,\lambda_N$ the eigenvalues of $K/N$. The Vendi Score is the exponential of
the von Neumann (matrix) entropy of that kernel,
\[
\mathrm{VS}(E)=\exp\Bigl(-\sum_i \lambda_i \log \lambda_i\Bigr),
\]
and is interpretable as the \emph{effective number of distinct questions}: a set of $N$
mutually orthogonal questions scores $N$, while $N$ paraphrases of one question score $1$.
Unlike APD, it is insensitive to duplicated modes and therefore penalises clustered
redundancy directly.

Figure~\ref{fig:semantic-dblp} shows the $t$-SNE projection of both question sets.
\textsc{QueryBloom} attains $\mathrm{APD}=0.385$ and $\mathrm{VS}=8.62$, against
$\mathrm{APD}=0.423$ and $\mathrm{VS}=10.99$ for the baseline. The two point clouds occupy
comparable regions of the embedding space with no visible collapse or clustering on either
side: on this axis the two benchmarks are essentially on par, with the human set retaining a
small edge. This is the expected --- and the intended --- outcome. Question embeddings are
dominated by lexical and topical variation, which human annotators produce effortlessly and
which is largely orthogonal to whether a benchmark actually stresses a query engine.
Semantic diversity is thus best read as a \emph{sanity floor}: it confirms that
\textsc{QueryBloom} does not generate paraphrase-like questions from a handful of prompts,
so that the gains reported on the remaining axes are not bought by sacrificing linguistic
variety. It is precisely because the input space is matched that the divergence in the
output and program spaces is attributable to the benchmarks themselves.

\subsection{Axis 2: Output (answer) diversity}

A question is characterised not only by how it is phrased but by the \emph{shape of its
answer}. We execute every query against the DBLP endpoint and record two per-query
quantities: the number of result \emph{rows} (answer cardinality) and the number of result
\emph{columns} (projected variables, i.e.\ the arity of the answer). Panel~(b) of
Figure~\ref{fig:diversity-dblp} shows the back-to-back distributions of both, with medians
marked.

The baseline is concentrated on a single answer shape: its median answer has $1$ row and
$1$ column (mean arity $1.02$, column standard deviation $0.47$), i.e.\ almost every
question asks for one atomic value or a boolean. \textsc{QueryBloom} spans a genuinely
multi-modal output space --- median $2$ rows and $3$ columns, with column standard deviation
$1.78$ and row standard deviation $36.7$ versus $16.8$ for the baseline. In other words, it
mixes single-fact lookups with list-valued, aggregate, and multi-attribute tabular answers.
This matters for evaluation, because systems that are only ever asked for one value are
never tested on projection, ordering, or grouping correctness, and the metrics used to score
them (exact match on a single cell) do not generalise to the tabular answers that real users
request.

\subsection{Axis 3: SPARQL (program) diversity}

The third axis characterises the queries themselves. Panel~(a) of
Figure~\ref{fig:diversity-dblp} summarises it as a radar plot; each spoke is normalised
independently by the larger of the two benchmark values, so the shape encodes relative
standing while the printed numbers give the raw values.

\paragraph{Template ratio.} Each query is reduced to a structural skeleton by stripping
\texttt{PREFIX} declarations and replacing every variable, IRI, literal, and prefixed name
by a placeholder; the metric is the number of distinct skeletons divided by $N$ (maximum
$1.0$). \textsc{QueryBloom} scores $1.000$: all $83$ queries are structurally distinct. The
baseline scores $0.627$, meaning that over a third of its queries are re-instantiations of a
skeleton already present --- the same graph pattern with different entities. The same
picture holds inside filter expressions, where the analogous \emph{filter template ratio}
is $0.338$ for \textsc{QueryBloom} against $0.114$ for the baseline: the baseline reuses a
small number of filter idioms almost exclusively.

\paragraph{Entity diversity.} The ratio of distinct entity IRIs to total entity mentions
across the benchmark. Here the baseline is nominally higher ($0.695$ vs.\ $0.601$). The
metric is normalised by \emph{mentions}, so it is mechanically depressed by queries that
reference several entities at once, and \textsc{QueryBloom} deliberately revisits the same
entities in different structural roles --- as a filter argument, as a join anchor, as a
grouping key --- to isolate the effect of query structure from the effect of topic. As the
coverage figures below show, this is not schema narrowness: \textsc{QueryBloom} touches
strictly more of the ontology than the baseline while reusing individual entities more
often.

\paragraph{Feature entropy.} Let $f_k$ be the fraction of queries using SPARQL feature $k$
over the $13$ tracked features, and $p_k=f_k/\sum_j f_j$. Feature entropy is the Shannon
entropy $-\sum_k p_k\log_2 p_k$, maximised at $\log_2 13 = 3.70$ when the benchmark's usage
mass is spread evenly across all features and minimised when it concentrates on a few.
\textsc{QueryBloom} reaches $3.241$ of the $3.70$ ceiling; the baseline reaches $2.697$,
reflecting a distribution dominated by plain \texttt{SELECT}/\texttt{ASK} patterns.

\paragraph{Complexity spread.} Panel~(b) additionally reports the number of structural
operators per query (all tracked features excluding the \texttt{SELECT}/\texttt{ASK} query
form). \textsc{QueryBloom} has a median of $3$ operators per query (mean $3.25$, standard
deviation $1.51$) against a median of $1$ for the baseline (mean $1.47$, standard deviation
$0.998$). The relevant observation is not merely that \textsc{QueryBloom} queries are
harder, but that they are harder \emph{at a range of difficulties}: the baseline
distribution is a spike at one operator, so it cannot separate systems by the complexity
they can sustain, whereas \textsc{QueryBloom}'s wider spread yields a graded difficulty
signal.

\paragraph{Feature distribution.} Panel~(c) breaks the aggregate entropy into per-feature
usage rates (percentage of queries containing each feature; \textsc{QueryBloom} above the
axis, baseline below). \textsc{QueryBloom} dominates on every construct that exercises
non-trivial query planning: \texttt{OPTIONAL} $27.7\%$ vs.\ $0.0\%$, \texttt{ORDER BY}
$57.8\%$ vs.\ $3.6\%$, \texttt{GROUP BY} $24.1\%$ vs.\ $3.6\%$, \texttt{LIMIT} $21.7\%$
vs.\ $3.6\%$, \texttt{FILTER} $54.2\%$ vs.\ $37.3\%$, and property paths $7.2\%$ vs.\
$0.0\%$; \texttt{UNION} is the one construct where the two are comparable ($7.2\%$ vs.\
$8.4\%$). Conversely, the baseline's mass concentrates on the \texttt{ASK} form
($34.9\%$ vs.\ $16.9\%$), whose answers are a single boolean. Three constructs
(\texttt{OPTIONAL}, \texttt{HAVING}, property paths) never appear in the baseline at all,
which means no system evaluated on it is ever tested on optional matching, post-aggregation
filtering, or path traversal.

\subsection{Comprehensiveness: how much is covered}

Diversity says how different the items are from one another; comprehensiveness says how much
of the target space they collectively occupy. We define three coverage ratios, all computed
against the DBLP ontology index and reported as raw fractions on the radar in
panel~(a) of Figure~\ref{fig:diversity-dblp}.

\paragraph{Class coverage.} For every query we resolve, from its AST, the ontology classes
it constrains --- both explicitly (via \texttt{rdf:type} or class IRIs appearing in triple
patterns) and implicitly, by resolving the types of the concrete entities it mentions
against the endpoint. Class coverage is the size of the union of these classes over the
benchmark, intersected with the ontology's class set to discard meta-level IRIs, divided by
the number of ontology classes. \textsc{QueryBloom} covers $12/29$ classes ($41.4\%$)
against $8/29$ ($27.6\%$) for the baseline.

\paragraph{Property coverage.} Defined analogously over predicate positions: the fraction of
ontology properties that appear in at least one query. This is the sharpest separation in
the study --- $35/81$ ($43.2\%$) for \textsc{QueryBloom} against $8/81$ ($9.9\%$) for the
baseline, a $4.4\times$ gap. The baseline exercises eight relations of the DBLP schema;
a system could score well on it while failing on nine tenths of the graph's vocabulary.

\paragraph{Feature coverage.} The fraction of the $13$ tracked SPARQL~1.1 constructs used at
least once. \textsc{QueryBloom} attains $13/13$ ($1.00$) --- full coverage of the tracked
language surface --- against $10/13$ ($0.77$) for the baseline.

\subsection{Summary}

Across the three diversity axes and the coverage axis, the pattern is consistent.
\textsc{QueryBloom} matches the human-curated baseline where input-side linguistic variety is
concerned (APD $0.385$ vs.\ $0.423$; Vendi $8.62$ vs.\ $10.99$), and improves on it
substantially everywhere the benchmark's discriminative power actually resides: structurally
unique queries ($1.000$ vs.\ $0.627$), balance of language-feature usage (entropy $3.241$
vs.\ $2.697$ out of $3.70$), operator complexity (median $3$ vs.\ $1$), answer-shape variety
(median arity $3$ vs.\ $1$), schema coverage ($41.4\%$ vs.\ $27.6\%$ of classes; $43.2\%$
vs.\ $9.9\%$ of properties), and language coverage ($13/13$ vs.\ $10/13$ features). The
matched semantic axis is what licenses this reading: the two benchmarks pose comparably
varied questions, and the difference lies in what those questions demand of the system that
must answer them.
